\section*{Bab \ref{ch:g11:enzymes-and-phenotype} --- Enzim dan Fenotipe}

\begin{solution}{exo:g11:enzymes-and-phenotype:1}
\omterm{def:g11:enzymes-and-phenotype:phenotype}{Genotipe}: perangkat \omterm{def:g10:universal-dna:gene}{alel} yang dibawanya. \omterm{def:g11:enzymes-and-phenotype:phenotype}{Fenotipe}: perangkat ciri yang
dapat diamati, pada skala molekuler, seluler dan organisme.
\end{solution}

\begin{solution}{exo:g11:enzymes-and-phenotype:2}
Kantong pada \omterm{def:g11:enzymes-and-phenotype:enzyme}{enzimnya} tempat \omterm{def:g11:enzymes-and-phenotype:enzyme}{substratnya} mengikat dan diubah.
Kekhususan: tapaknya cocok dengan satu \omterm{def:g11:enzymes-and-phenotype:enzyme}{substrat} (atau satu keluarga
\omterm{def:g11:enzymes-and-phenotype:enzyme}{substrat}) dan mengatalisis satu reaksi.
\end{solution}

\begin{solution}{exo:g11:enzymes-and-phenotype:3}
Mendekati \qty{37}{\celsius}; pada \qty{45}{\celsius} sekitar sepertiganya tersisa.
\end{solution}

\begin{solution}{exo:g11:enzymes-and-phenotype:4}
Molekuler: tirosinasenya tidak ada atau tidak aktif, tidak ada melanin
yang terbuat. Seluler: \omterm{def:g10:cells-common-unit:cell}{sel} pigmennya hadir tetapi kosong dari melanin.
Organisme: rambut putih, kulit dan mata yang sangat pucat, kepekaan
terhadap cahaya dan mudah terbakar matahari.
\end{solution}

\begin{solution}{exo:g11:enzymes-and-phenotype:5}
Lipatan pepsin hanya bekerja dalam asam yang kuat (pH 2); ususnya
netral sampai sedikit \omterm{def:g10:universal-dna:nucleotide}{basa}, tempat lengkungnya menunjukkan kegiatan nol.
\end{solution}

\begin{solution}{exo:g11:enzymes-and-phenotype:6}
Satu \omterm{def:g10:universal-dna:gene}{alel} yang bekerja menghasilkan separuh \omterm{def:g11:enzymes-and-phenotype:enzyme}{enzim} yang normal, tetapi
satu molekul \omterm{def:g11:enzymes-and-phenotype:enzyme}{enzim} bekerja ribuan kali tiap sekon: separuh jumlahnya
tetap mengubah cukup tirosin untuk membuat pigmen yang normal.
\omterm{def:g11:enzymes-and-phenotype:phenotype}{Fenotipenya} sudah jenuh; hanya dua \omterm{def:g10:universal-dna:gene}{alel} yang gagal yang terlihat.
\end{solution}

\begin{solution}{exo:g11:enzymes-and-phenotype:7}
Albinismenya berasal dari halangan pada dua \omterm{def:g10:universal-dna:gene}{gen} yang berbeda (dua \omterm{def:g11:enzymes-and-phenotype:enzyme}{enzim}
jalurnya yang berbeda, atau pengangkutan melaninnya). Anaknya mewarisi
satu \omterm{def:g10:universal-dna:gene}{alel} yang bekerja bagi tiap \omterm{def:g10:universal-dna:gene}{gennya} dari orang tua yang halangannya
di tempat lain, sehingga kedua \omterm{def:g11:enzymes-and-phenotype:enzyme}{enzimnya} bekerja dan pigmennya terbuat.
\end{solution}

\begin{solution}{exo:g11:enzymes-and-phenotype:8}
Fenilalaninnya menumpuk; tirosinnya (dan apa yang berasal darinya)
kurang. Karena melanin dibuat dari tirosin, kekurangannya mengurangi
pigmennya: rambut dan kulit yang terang lazim pada fenilketonuria yang
tidak diobati.
\end{solution}

\begin{solution}{exo:g11:enzymes-and-phenotype:9}
Di tempat susu menjadi makanan pokok, mencernanya menyediakan energi
dan kalsium tanpa membuat sakit, sehingga \omterm{def:g10:universal-dna:gene}{alel} yang lestari itu menjadi
keuntungan. Di tempat orang dewasanya tidak minum susu, \omterm{def:g10:universal-dna:gene}{alelnya} tidak
mengubah apa pun; ia ``lebih sehat'' hanya dalam lingkungan yang
memasok \omterm{def:g11:enzymes-and-phenotype:enzyme}{substratnya}.
\end{solution}

\begin{solution}{exo:g11:enzymes-and-phenotype:10}
Untuk tirosinase: molekuler, \omterm{def:g11:enzymes-and-phenotype:enzyme}{substratnya} tidak lagi cocok, tidak ada
melanin; seluler, \omterm{def:g10:cells-common-unit:cell}{sel} pigmennya kosong; organisme, albinisme. Untuk
laktase: laktosanya tidak tercerna; \omterm{def:g10:cells-common-unit:cell}{sel} ususnya kekurangan gula itu dan
usus besarnya memfermentasinya; ketidaktahanan.
\end{solution}

\begin{solution}{exo:g11:enzymes-and-phenotype:11}
Ujung tubuhnya tetap lebih hangat daripada \qty{35}{\celsius},
\omterm{def:g11:enzymes-and-phenotype:enzyme}{enzimnya} nyaris tidak bekerja, dan anak kucingnya pucat di seluruh
tubuhnya. Di luar rumah pada musim dingin telinga dan telapaknya
mendingin, \omterm{def:g11:enzymes-and-phenotype:enzyme}{enzimnya} menjadi aktif, dan bulu yang tumbuh kembali di sana
muncul gelap.
\end{solution}

\begin{solution}{exo:g11:enzymes-and-phenotype:12}
Molekuler: separuh hemoglobinnya berupa bentuk sabitnya. Seluler: pada
oksigen yang normal \omterm{def:g10:cells-common-unit:cell}{selnya} tetap bulat; pada oksigen yang sangat rendah
sebagian menyabit. Organisme: sehat dalam kehidupan biasa, berisiko di
tempat tinggi atau ketika menyelam --- dan terlindung dari malaria.
Lingkungannya (oksigen, malaria) yang memutuskan apakah \omterm{def:g11:enzymes-and-phenotype:phenotype}{fenotipe}
pembawanya menjadi beban atau keuntungan.
\end{solution}

\begin{solution}{exo:g11:enzymes-and-phenotype:13}
Pada \qty{41}{\celsius} \omterm{def:g11:enzymes-and-phenotype:enzyme}{enzimnya} masih menahan sebagian besar
kegiatannya (lengkungnya dekat ke puncaknya); pada \qty{43}{\celsius}
kegiatannya turun tajam seiring \omterm{def:g11:gene-expression:protein}{proteinnya} yang mulai terbuka
lipatannya, dan jika terbukanya berlanjut, ia tak terbalikkan: \omterm{def:g11:enzymes-and-phenotype:enzyme}{enzimnya} musnah
alih-alih melambat.
\end{solution}

\begin{solution}{exo:g11:enzymes-and-phenotype:14}
Lebih banyak \omterm{def:g10:cells-common-unit:organelle}{mitokondria} dan kapiler tiap \omterm{def:g10:muscles-and-joints:muscle}{serat ototnya}; jantung yang
lebih besar dengan volume sekuncup yang lebih besar; lebih banyak
hemoglobin tiap liter darahnya; simpanan \omterm{def:g10:exercise-and-energy:glycogen}{glikogen} yang lebih besar.
Masing-masing adalah \omterm{def:g11:enzymes-and-phenotype:phenotype}{fenotipe} yang dihasilkan \omterm{def:g11:enzymes-and-phenotype:phenotype}{genotipe} yang sama di
bawah lingkungan yang berbeda: sifat ``genetik'' adalah rentang yang
dibolehkan \omterm{def:g11:enzymes-and-phenotype:phenotype}{genotipenya}, bukan sebuah nilai yang tetap.
\end{solution}

\begin{solution}{exo:g11:enzymes-and-phenotype:15}
\omterm{def:g11:enzymes-and-phenotype:phenotype}{Genotipenya} menetapkan \omterm{def:g11:enzymes-and-phenotype:enzyme}{enzim} mana yang dibuat dan bagaimana ia
berperilaku: tanpa \omterm{def:g11:enzymes-and-phenotype:enzyme}{enzim} 1 (fenilketonuria), tanpa tirosinase
(albinisme), tirosinase yang peka panas (Siam). Lingkungannya dapat
mengubah hasilnya hanya dalam batas yang dibolehkan \omterm{def:g11:enzymes-and-phenotype:enzyme}{enzimnya}: pola
makan menyingkirkan \omterm{def:g11:enzymes-and-phenotype:enzyme}{substratnya} dan menyelamatkan otaknya, tetapi tidak
dapat membuat \omterm{def:g11:enzymes-and-phenotype:enzyme}{enzim} yang hilang; dingin membuat \omterm{def:g11:enzymes-and-phenotype:enzyme}{enzim} Siamnya bekerja,
tetapi tidak ada suhu yang membuat \omterm{def:g11:enzymes-and-phenotype:enzyme}{enzim} albinonya bekerja; dan tidak
ada apa pun dalam lingkungannya yang memberi si albino pigmennya.
\omterm{def:g11:enzymes-and-phenotype:phenotype}{Genotipenya} menetapkan rentangnya; lingkungannya memilih di dalamnya.
\end{solution}

\begin{solution}{pb:g11:enzymes-and-phenotype:1}
\textbf{1.} Badan \qty{38}{\celsius}: sekitar 2\%, pucat. Kaki
\qty{36}{\celsius}: sekitar 20\%, pucat. Telapak \qty{33}{\celsius}:
80\%, gelap. Telinga dan moncong \qty{31}{\celsius}: lebih dari 90\%, gelap.

\textbf{2.} Badan dan kaki yang pucat, telapak, telinga, moncong (dan
ekor) yang gelap: itulah pola Siamnya.

\textbf{3.} 100\% di mana-mana: seekor kucing yang berpigmen merata.

\textbf{4.} Dalam rahimnya tiap bagian anak kucingnya bersuhu
\qty{38}{\celsius}, \omterm{def:g11:enzymes-and-phenotype:enzyme}{enzimnya} tidak aktif di seluruh tubuhnya, dan tidak
ada pigmen yang tertimbun sebelum lahir; titiknya menggelap seiring
ujung tubuhnya mendingin setelah lahir.

\textbf{5.} Sekitar \qty{35.5}{\celsius}, tempat kegiatannya melintasi
30\%: bagian lengkungnya yang menurun tajam antara 33 dan
\qty{37}{\celsius} yang menetapkan batasnya.

\textbf{6.} Gelap: pada \qty{30}{\celsius} \omterm{def:g11:enzymes-and-phenotype:enzyme}{enzimnya} aktif sepenuhnya
dalam bulu yang sedang tumbuh.

\textbf{7.} Pucat: pada \qty{38}{\celsius} \omterm{def:g11:enzymes-and-phenotype:enzyme}{enzimnya} tidak aktif.

\textbf{8.} Pigmennya ditimbun dalam bulunya sementara bulunya tumbuh
dan tidak dapat disingkirkan atau ditambahkan sesudahnya; warnanya
berubah hanya ketika bulunya rontok dan digantikan bulu baru yang
tumbuh pada suhu setempat.

\textbf{9.} \omterm{def:g11:enzymes-and-phenotype:phenotype}{Genotipe} \omterm{def:g10:cells-common-unit:cell}{sel} bidang itu sama sebelum dan sesudahnya, dan
sama dengan sisa badannya; hanya suhu \omterm{def:g11:enzymes-and-phenotype:enzyme}{enzimnya} yang berubah, dan
bersamanya pigmennya. Tindakannya mengenai kegiatan \omterm{def:g11:enzymes-and-phenotype:enzyme}{enzimnya}, yaitu
peristiwa molekuler.

\textbf{10.} Seluruh \omterm{def:g10:cells-common-unit:cell}{sel} kucingnya berasal dari satu \omterm{def:g10:cells-common-unit:cell}{sel} telur dan
membawa \omterm{def:g10:universal-dna:gene}{alel} yang sama; dan percobaan bidang cukurnya mengubah warna
sebuah daerah tanpa mengubah \omterm{def:g10:cells-common-unit:cell}{selnya} --- sebuah \omterm{def:g10:universal-dna:gene}{alel} tidak dapat
dialihkan oleh kantong es.

\textbf{11.} Molekuler: tirosinase yang hanya bekerja di bawah sekitar
\qty{35}{\celsius}. Seluler: \omterm{def:g10:cells-common-unit:cell}{sel} pigmen yang penuh melanin pada ujung
tubuhnya, kosong pada badannya. Organisme: kucing pucat dengan titik
yang gelap.

\textbf{12.} \omterm{def:g10:chemistry-of-life:families}{Asam aminonya} duduk di tempat ia memantapkan lipatannya;
penggantinya menahan lipatannya tetap bersatu pada suhu yang rendah
tetapi membiarkannya mengendur beberapa derajat lebih tinggi, sehingga
\omterm{def:g11:enzymes-and-phenotype:enzyme}{tapak aktifnya} kehilangan bentuknya lebih dahulu daripada \omterm{def:g11:enzymes-and-phenotype:enzyme}{enzim} biasanya.

\textbf{13.} \omterm{def:g11:enzymes-and-phenotype:enzyme}{Enzim} \omterm{def:g10:universal-dna:gene}{alel} biasanya aktif sepenuhnya pada suhu tubuhnya;
separuh jumlah yang normal masih jauh lebih dari cukup untuk membuat
pigmen di mana-mana. \omterm{def:g10:universal-dna:gene}{Alel} Siamnya bersifat resesif terhadapnya.

\textbf{14.} Molekuler: tidak ada \omterm{def:g11:enzymes-and-phenotype:enzyme}{enzim} yang bekerja pada suhu berapa
pun, berhadapan dengan \omterm{def:g11:enzymes-and-phenotype:enzyme}{enzim} yang bekerja bersyarat. Seluler: tidak ada
melanin di mana pun, berhadapan dengan melanin pada daerah yang sejuk.
Organisme: kucing putih seluruhnya dengan mata merah muda, berhadapan
dengan pola bertitiknya. Dingin tidak mengubah apa pun bagi si albino.

\textbf{15.} Bahwa keduanya membawa \omterm{def:g10:universal-dna:gene}{alel} \omterm{def:g10:universal-dna:gene}{gen} yang sama yang peka suhu;
barangkali substitusi serupa yang menggoyahkan \omterm{def:g11:enzymes-and-phenotype:enzyme}{enzimnya}, yang muncul
sendiri-sendiri pada tiap \omterm{def:g10:biodiversity-scales:species}{spesiesnya}.

\textbf{16.} Dua puluh kali lipat; \omterm{def:g11:enzymes-and-phenotype:enzyme}{enzim} 1, yang mengubah fenilalanin
menjadi tirosin.

\textbf{17.} Fenilalanin yang berlebih mengganggu \omterm{def:g10:cells-common-unit:cell}{sel} otak yang sedang
berkembang secara khusus (pasokan \omterm{def:g10:chemistry-of-life:families}{asam amino} lainnya dan pematangannya);
\omterm{def:g10:cells-common-unit:cell}{sel} kulitnya menahannya. Cacat molekulernya sama di mana-mana; akibat
selulernya tidak.

\textbf{18.} Seperlima dari \num{1200} adalah \qty{240}{\micro mol/L},
di bawah \num{360}: aman. Tanpa \omterm{def:g11:enzymes-and-phenotype:enzyme}{enzimnya}, aras darahnya sekadar
mengikuti asupannya, sehingga \omterm{def:g11:enzymes-and-phenotype:phenotype}{fenotipenya} ditetapkan pasokan
\omterm{def:g11:enzymes-and-phenotype:enzyme}{substratnya} --- yaitu lingkungannya --- alih-alih oleh \omterm{def:g11:enzymes-and-phenotype:enzyme}{enzim} yang
hilang itu sendirian.

\textbf{19.} Fenilalanin diperlukan untuk membangun \omterm{def:g11:gene-expression:protein}{protein} tubuhnya
dan tidak dapat disusun sendiri; tanpa sama sekali, pertumbuhannya
berhenti. Asupannya harus cukup bagi penyusunan \omterm{def:g11:gene-expression:protein}{proteinnya} dan tidak
lebih, dan itu menuntut pengukuran serta penyetelan.

\textbf{20.} \omterm{def:g11:enzymes-and-phenotype:phenotype}{Genotipenya} menetapkan \omterm{def:g11:enzymes-and-phenotype:enzyme}{enzimnya} (peka panas pada
kucingnya, tidak ada pada anaknya); lingkungannya menetapkan syaratnya
(suhu kulit; fenilalanin dalam pola makannya); \omterm{def:g11:enzymes-and-phenotype:phenotype}{fenotipenya} adalah hasil
\omterm{def:g11:enzymes-and-phenotype:phenotype}{genotipe} yang terungkap dalam sebuah lingkungan tertentu.
\end{solution}
